\exercisesfor{9}

\begin{enumerate}
\def\labelenumi{\arabic{enumi}.}
\item
  令 \(\mathrm{G} = \mathbf{SO}(3, \mathbf{R})\). 令 \(\Delta_{1}, \Delta_{2}\) 为 \(R^{3}\) 中的两条正交直线, 令 \(H_{i}\) 为 G 中绕 \(\Delta_{i} (i = 1, 2)\) 的旋转所成的子群 . 则 \([\mathrm{L}(\mathrm{H}_{1}), \mathrm{L}(\mathrm{H}_{2})]\) 是 G 的一个 1 维李子代数, 不同于 \((\mathrm{H}_{1}, \mathrm{H}_{2})\) 在 e 处相切的李子代数.
\item
  假设 K 是超度量域. 设 G 是有限维李群, g 是其李代数, A 是 G 的有限子集. 则 \(Z_{G}(A)\) 是 G 的李子群, 其李代数为 \(\mathfrak{z}_{\mathfrak{g}}(A)\). (仿照命题 8.)
\item
  设 G 是连通实或复李群. G 的中心 Z 是 G 的李拟子群, 且 L(Z) 是 L(G) 的中心.
\item
  假设 K 是超度量域且 p \textgreater{} 0 (采用 §7 的记号). 设 G 是有限维李群, A 是 G 的自同构群, B 是相应的 L(G) 的自同构群. 令 \(\mathrm{G}^{\mathbf{A}}\) (分别为 \(\mathrm{L}(\mathrm{G})^{\mathbf{B}}\)) 表示 G (分别为 L(G)) 中在 A (分别为 B) 下不动的元素的集合. 则 \(G^{A}\) 是 G 的李子群, 其李代数为 \(\mathrm{L}(G)^{\mathbf{B}}\). (使用对数映射.)
\item
  设 \(G\) 是有限维连通实或复李群. 令 \((G_0, G_1, \ldots)\) 为 (第 II 章 § 4 练习 18 的) 上中心列, 令 \((g_0, g_1, \ldots)\) 为 L(G) 的上中心列 (第 I 章 § 1 第 6 号). 则对所有 \(L(G)\)\(i\), \(G_i\) 是 G 的李子群, 且 \(G\)\(L(G_i) = g_i\).
\item
  令 \(\Gamma\) 为 §4 练习 5 (b) 中定义的 3 维单连通幂零实李群. 令 \(\alpha \in \mathbf{R}\) 为无理数. 令 P 为 \(P\)\(\Gamma \times \mathbf{R}\) 的离散子群, 由 \(((0,0,x),\alpha x)\) (\(x \in \mathbf{Z}\)) 组成. 令 \(G = (\Gamma \times \mathbf{R}) / P\). 证明 \((G,G)\) 在 G 中不是闭的.\(G\)
\item
  \begin{enumerate}
  \def\labelenumii{(\alph{enumii})}
  \tightlist
  \item
    令 \(G\) 为连通半单实李群, Z 为其中心, \(Z\)\(\rho\) 为 G 在有限维复向量空间上的连续线性表示. 存在整数 p, 使得对所有 \(G\)\(p\)\(z \in Z\), \(\rho(z)\) 都可对角化, 且 \(\rho(z)\) 的所有特征值都是 p 次单位根. (可假设 \(p\)\(\rho\) 不可约. 根据 Schur 引理, \(\rho(z)\) 是标量. 另一方面, 对所有  有 det \(\rho(g) = 1\), 因为 \(g \in G\)\(G = \mathcal{D}G\).)
  \end{enumerate}
\end{enumerate}

\begin{enumerate}
\def\labelenumi{(\alph{enumi})}
\setcounter{enumi}{1}
\tightlist
\item
  推出若 G 具有忠实的有限维连续线性表示, 则 Z 是有限的.\(G\)\(Z\)
\end{enumerate}

¶ 8. 设 G 是有限维连通实李群, g 是其李代数, n 是 g 的最大幂零理想, 且 h = {[}g, g{]} + n.

\begin{enumerate}
\def\labelenumi{(\alph{enumi})}
\item
  \(\mathfrak{h}\) 是 \(\mathfrak{g}\) 的特征理想; \(\mathfrak{h}\) 的根为 n; 对所有 \(n\)\(x \in \mathfrak{h}\), \(\operatorname{Tr} \operatorname{ad}_{\mathfrak{h}} x = 0\); 对于 \(l\)\(\mathfrak{g}\) 的每个 Levi 截面 l, \(\mathfrak{h} = l + n\).
\item
  假设 G 是单连通的. 令 Z 为其中心, H 为李代数为 b 的 G 的李子群, \(\phi\) 为从 G 到 G/H 的典范态射. 则 \(\phi(Z)\) 在 G/H 中是离散的. (群 G 是半单李群 S 与其根 R 的半直积. 令 \(z \in Z\). 则 \(z = y^{-1}x\), 其中 \(x \in R\), y 属于 S 的中心. 存在整数 p, 使 \(Ad_{q}y\) 的特征值, 从而 \(Ad_{q}x\) 的特征值也是 p 次单位根 (练习 7). 推出若 N 是李代数为 n 的 G 的李子群, 则 R 中存在 e 的一个邻域 U, 满足 U = UN = NU,
\end{enumerate}

\[
\mathrm{Z} \cap (\mathrm{SU}) \subset \mathrm{SN} = \mathrm{H}.
\]

\begin{enumerate}
\def\labelenumi{(\alph{enumi})}
\setcounter{enumi}{2}
\item
  现在不再假设 G 是单连通的. 令 H 为李代数为 b 的 G 的积分子群. 证明 H 是 G 的幺模正规李子群, 具有幂零根, 且 G/H 是交换的. (使用 (a) 和 (b).)
\item
  推出若  在 G 中稠密, 则 G 的根是幂零的.
\end{enumerate}

\begin{enumerate}
\def\labelenumi{\arabic{enumi}.}
\setcounter{enumi}{8}
\tightlist
\item
  令 \(G\) 为 \(\mathbf{SL}(2, \mathbf{R})\) 的万有覆盖. 将典范态射 \(G \to \mathbf{SL}(2, \mathbf{R})\) 的核与 \(\mathbf{Z}\) 识别 (§ 6, 练习 2). 令 a 为 \(a\)\(\mathbf{T}^n\) 中幂在 \(\mathbf{T}^n\) 中处处稠密的元素 (《一般拓扑》第 VII 章 §1 命题 7 的推论 2). 令 D 为 \(D\)\(G \times \mathbf{T}^n\) 中由 (1, a) 生成的离散子群. 令 \(H = (G \times \mathbf{T}^n)/D\). 则
\end{enumerate}

\[
\mathrm{L} (\mathrm{H}) = \mathfrak {s l} (2, \mathbf {R}) \times \mathbf {R} ^ {n}
\]

且李代数为  的 H 的积分子群 \(\mathbf{H}'\) 同构于 \(\mathbf{H}\)\(\mathfrak{sl}(2,\mathbf{R})\)\(\mathbf{G}\), 并在 \(\mathbf{H}\) 中稠密. \(\mathrm{D}^n\mathrm{H}' = \mathrm{H}'\) 对所有 \(n\geqslant 0\) 成立.

¶ 10. 设 G 是有限维实或复李群. 令

\[
p = \dim [ L (G), L (G) ].
\]

令 \((\mathrm{G},\mathrm{G})\) 具有其作为 \(\mathbf{G}\) 的积分子群的结构. 存在 \(\mathbf{V}\)\(e\)\((\mathrm{G},\mathrm{G})\) 中 e 的一个邻域 \(\mathbf{V}\), 使得 \(p\) 中的每个元素都是 \(\mathbf{G}\) 中元素的 p 个换位子的乘积. (令 \(x_{1},y_{1},\ldots ,x_{p},y_{p}\) 为 \(\mathbf{L}(\mathbf{G})\) 中的元素, 使 \([x_1,y_1]\) 构成 \(\mathbf{L}(\mathbf{G})\) 的一个基. 对于  中的 \(s,t\), 记\(\mathbf{R}\)

\[
\rho_ {i} (s, t) = (\exp s y _ {i}) ^ {- 1} (\exp t x _ {i}) ^ {- 1} (\exp s y _ {i}) (\exp t x _ {i}).
\]

将隐函数定理应用于映射

\[
\left(s _ {1}, t _ {1}, \dots , s _ {p}, t _ {p}\right) \mapsto \rho_ {1} \left(s _ {1}, t _ {1}\right) \dots \rho_ {p} \left(s _ {p}, t _ {p}\right)
\]

从 \(\mathbf{R}^{2p}\) 到 (G, G).

\begin{enumerate}
\def\labelenumi{\arabic{enumi}.}
\setcounter{enumi}{10}
\item
  设 G 是 B = Z/2Z 与 K 按李群 K 的自同构 \(x \mapsto -x\) 对应的半直积. 则 L(K) 是 L(G) 的理想, L(B) = \{0\}, 但 (K, B) = K.
\item
  设 \((e_1, e_2, e_3)\) 是 \(\mathbf{K}^3\) 的典范基. 考虑 \(\mathbf{K}^3\) 上的幂零李代数结构, 使得 \([e_1, e_2] = e_3\), \([e_1, e_3] = [e_2, e_3] = 0\). 赋予 \(\mathbf{K}^3\) 相应的群法则 (第 5 号),
\end{enumerate}

\[
(x, y, z) (x ^ {\prime}, y ^ {\prime}, z ^ {\prime}) = (x + x ^ {\prime}, y + y ^ {\prime}, z + z ^ {\prime} + \frac {1}{2} (x y ^ {\prime} - y x ^ {\prime})).
\]

映射

\[
(\lambda , \mu , \nu) \mapsto (\exp \lambda e _ {1}) (\exp \mu e _ {2}) (\exp \nu (e _ {1} + e _ {3}))
\]

从 \(K^{3}\) 到 G 的映射既不是满射 (证明 \((0,1,1)\) 不在其像中), 也不是单射 (证明 \((0,1,0)\) 与 \((1,1,-1)\) 的像相同).

\begin{enumerate}
\def\labelenumi{\arabic{enumi}.}
\setcounter{enumi}{12}
\tightlist
\item
  \begin{enumerate}
  \def\labelenumii{(\alph{enumii})}
  \tightlist
  \item
    在 \(\mathbf{R}^3\) 上定义如下乘法:
  \end{enumerate}
\end{enumerate}

\[
(x, y, z). (x ^ {\prime}, y ^ {\prime}, z ^ {\prime}) = (x + x ^ {\prime} \cos z - y ^ {\prime} \sin z, y + x ^ {\prime} \sin z + y ^ {\prime} \cos z, z + z ^ {\prime}).
\]

证明由此得到一个可解实李群 \(\mathbf{G}\), 使得 \(\mathrm{DG} = \mathbf{R}^2\times \{0\}\).  的中心 \(\mathbf{Z}\) 为 \(\mathbf{G}\)\(\{0\} \times \{0\} \times 2\pi \mathbf{Z}\).

\begin{enumerate}
\def\labelenumi{(\alph{enumi})}
\setcounter{enumi}{1}
\tightlist
\item
  对于 \((x,y,z)\in G\), 令 \(\pi (x,y,z)\) 为映射
\end{enumerate}

\[
(\lambda , \mu) \mapsto (\lambda \cos z - \mu \sin z + x, \lambda \sin z + \mu \cos z + y)
\]

从 \(\mathbf{R}^2\) 到 \(\mathbf{R}^2\) 的映射. 证明 \(\pi\) 是从 G 到仿射群 \(G\) 的一个李子群 \(\mathbf{G}'\) 的态射, 该子群由 \(\mathbf{R}^2\) 的平移和旋转生成. 证明 \(\mathbf{R}^2\)\(\operatorname{Ker} \pi = \mathbf{Z}\).

\begin{enumerate}
\def\labelenumi{(\alph{enumi})}
\setcounter{enumi}{2}
\item
  我们典范地将 \(\mathrm{L}(\mathrm{G}) = \mathrm{T}_{(0,0,0)}\mathrm{G}\) 与 \(\mathbf{R}^3\) 识别. 令 \((e_1, e_2, e_3)\) 为 \(\mathbf{R}^3\) 的典范基. 证明 \([e_1, e_2] = 0\), \([e_3, e_1] = e_2\), \([e_3, e_2] = -e_1\).
\item
  证明当 \(c \neq 0\) 时,
\end{enumerate}

\[
\exp_ {\mathrm{G}} (a, b, c) = \left(\frac {1}{c} (a \sin c + b \cos c - b), \frac {1}{c} (- a \cos c + b \sin c + a), c\right).
\]

推出对所有 \(u \in L(G)\), 都有 \(Z \subset \exp(\mathbf{R}u)\). 证明 \(\exp_G: L(G) \to G\) 既不是单射也不是满射.

\begin{enumerate}
\def\labelenumi{\arabic{enumi}.}
\setcounter{enumi}{13}
\tightlist
\item
  \begin{enumerate}
  \def\labelenumii{(\alph{enumii})}
  \tightlist
  \item
    令 \(\mathbf{G} = \mathbf{GL}(n, \mathbf{C})\). 证明 \(\exp_{\mathbf{G}}\) 是满射. (使用《谱理论》第 I 章 § 4 第 8 号的全纯函数演算.)
  \end{enumerate}
\end{enumerate}

\begin{enumerate}
\def\labelenumi{(\alph{enumi})}
\setcounter{enumi}{1}
\tightlist
\item
  令 \(\mathbf{G}' = \mathbf{SL}(2, \mathbf{C})\). 证明若 \(\sigma \in \mathbf{C}^*\), 则元素 \(\left( \begin{array}{cc} -1 & \sigma \\ 0 & -1 \end{array} \right)\)
\end{enumerate}

\(G'\) 不属于 \(\exp_{G'}\) 的像.

\begin{enumerate}
\def\labelenumi{\arabic{enumi}.}
\setcounter{enumi}{14}
\tightlist
\item
  \begin{enumerate}
  \def\labelenumii{(\alph{enumii})}
  \tightlist
  \item
    令 \(x \in \mathfrak{sl}(2, \mathbf{R})\), \(\Delta = \det x\). 证明
  \end{enumerate}
\end{enumerate}

\[
e ^ {x} = \operatorname{ch} \sqrt {- \Delta}. I + \frac {\operatorname{sh} \sqrt {- \Delta}}{\sqrt {- \Delta}} \cdot x \quad \mathrm{if} \Delta <   0
\]

\[
e ^ {x} = \cos {\sqrt {\Delta}}. I + \frac {\sin {\sqrt {\Delta}}}{\sqrt {\Delta}}. x \qquad \mathrm{if} \Delta > 0.
\]

\begin{enumerate}
\def\labelenumi{(\alph{enumi})}
\setcounter{enumi}{1}
\tightlist
\item
  令 \(\mathbf{H} = \mathbf{SL}(2, \mathbf{R})\). 证明 \(g = \begin{pmatrix} \lambda & 0 \\ 0 & \lambda^{-1} \end{pmatrix} \in \mathbf{H}\) 是 \(\exp_{\mathbf{H}}\) 的像, 当且仅当 \(\lambda > 0\) 或 \(\lambda = -1\). 若 \(\lambda = 0\), 则所有包含 g 的单参数子群都相等.\(g\)
\end{enumerate}

\begin{enumerate}
\def\labelenumi{\arabic{enumi}.}
\setcounter{enumi}{15}
\tightlist
\item
  设 G 是有限维李群. 证明存在 L(G) 的一个基 \((x_{1},\ldots,x_{n})\), 使对所有 i, \(\exp(\mathbf{R}x_{i})\) 都是 G 的李子群. (使用《谱理论》第 II 章 §2 引理 1.)
\end{enumerate}

¶ 17. (a) 令 c 表示具有基 \((a, b, c)\) 的实可解李代数, 满足 \([a, b] = c\), \([a, c] = -b\), \([b, c] = 0\). 令 d 表示具有基 \((a, b, c, d)\) 的实可解李代数, 满足 \([a, b] = c\), \([a, c] = -b\), \([b, c] = d\), \([a, d] = [b, d] = [c, d] = 0\). 令 g 为实可解李代数. 若 g 含有非零元素 x, y, z, 满足 \([x, y] = z\) 且 \([x, z] = -y\), 则 g 含有同构于 c 或 d 的子代数.~(令 \(a_{1} = y\), \(b_{1} = z\), \(c_{1} = [y, z]\); 递归地由 \(a_{i}\)\(b_{i}\)\(c_{i}\)\(a_{i} = [a_{i-1}, c_{i-1}]\), \(b_{i} = [b_{i-1}, c_{i-1}]\), \(c_{i} = [a_{i}, b_{i}]\) 定义 . 考虑满足 \(c_{k} = 0\) 的最小整数 k.)

\begin{enumerate}
\def\labelenumi{(\alph{enumi})}
\setcounter{enumi}{1}
\item
  设 g 和 \(g^{*}\) 是实可解李代数, \(\phi\) 是从 g 到 \(g^{*}\) 的满射同态. 若 \(g^{*}\) 含有同构于 c 或 d 的子代数, 则 g 也具有同样性质.
\item
  设 \(\mathfrak{h}\) 是复可解李代数. 考虑 \(\mathfrak{h}\) 的伴随表示的 Jordan-Hölder 列; 该列的商定义了 \(\mathfrak{h}\) 的一维表示, 从而定义了 \(\mathfrak{h}\) 上的线性形式. 这些只依赖于 \(\mathfrak{h}\) 的线性形式称为 \(\mathfrak{h}\) 的根. 若 \(\mathfrak{h}'\) 是实可解李代数, 则 \(\mathfrak{h}'\) 的根是  的根在 \(\mathfrak{h}'\) 上的限制.\(\mathfrak{h}' \otimes_{\mathbf{R}} \mathbf{C}\)
\end{enumerate}

然后令 G 为有限维单连通可解实李群, g 为其李代数. 证明以下条件等价: (α) exp₄ 是单射; (β) exp₅ 是满射; (γ) exp₆ 是双射; (δ) exp₇ 是从解析流形 L(G) 到解析流形 G 的同构; (ε) L(G) 不含同构于 c 或 d 的子代数; (ζ) L(G) 没有商代数含有同构于 c 的子代数; (η) g 的每个根都具有 \(\phi + i\delta'\) 的形式, 其中 \(\phi\), \(\phi'\) 属于 \(g^*\), 且 \(\phi'\) 与 \(\phi\) 成比例; (θ) 对所有 \(x \in g\),  的唯一纯虚特征值 (在 C 中) 是 0.†
% semantic-fragments: legacy-input-seam
\relax{}

\begin{enumerate}
\def\labelenumi{\arabic{enumi}.}
\setcounter{enumi}{17}
\tightlist
\item
  设 \(G\) 是连通实李群，\(Z\) 是其中心，\(Z_0\) 是 \(Z\) 的单位元连通分支，\(g\) 是 \(G\) 的李代数，而 \(h\) 是 \(g\) 的中心。于是 \(Z\) 和 \(Z_0\) 是 \(G\) 的李拟子群，其李代数为 \(h\)（练习 3）。\(g\) 上的可接受范数，是指定义 \(g\) 的拓扑并使 \(g\) 成为赋范李代数的范数。
\end{enumerate}

\begin{enumerate}
\def\labelenumi{(\alph{enumi})}
\item
  对所有 \(r > 0\) 及 \(z \in \mathfrak{z}\)，存在一个 \(\mathfrak{g}\) 上的可接受范数，其在 \(z\) 处的取值为 \(< r\)。（设 \(q\) 是 \(\mathfrak{g}\) 上的可接受范数。证明对于适当选取的 \(\lambda > 0\)，函数 \(x \mapsto \|x\| = \lambda q(x) + \inf_{y \in \mathfrak{z}} q(x + y)\) 具有所需性质。）
\item
  证明下列条件等价：
\item
  \(Z_{0}\) 是单连通的；
\end{enumerate}

\begin{enumerate}
\def\labelenumi{(\roman{enumi})}
\setcounter{enumi}{1}
\item
  对 g 上的每个可接受范数，\(\exp_{G}\) 在以 0 为中心、半径为 \(\pi\) 的开球上的限制是单射；
\item
  存在 \(r > 0\)，使得对于 \(\mathfrak{g}\) 上的每个可接受范数，\(\exp_{\mathbb{G}}\) 在以 0 为中心、半径为 \(r\) 的开球上的限制是单射。
\end{enumerate}

（为证明 (iii) \(\Rightarrow\) (i)，使用 (a)。为证明 (i) \(\Rightarrow\) (ii)，设 \(x, y\) 属于 \(\mathfrak{g}\)，\(\| x \| < \pi\)，\(\| y \| < \pi\)，\(x \neq y\)，且 \(\exp_{\mathfrak{g}} x = \exp_{\mathfrak{g}} y\)。于是 \(\exp \mathrm{d} x = \exp \mathrm{d} y\)，从而得到 \(\mathrm{ad} x = \mathrm{ad} y\)（将 §6 的命题 17 应用于 \(\mathfrak{g}\) 的复化）。因此存在一个非零 \(z\) 属于 \(\mathfrak{z}\)，使得 \(\exp_{\mathfrak{g}} z = e\)；于是 (i) 为假。）

\begin{enumerate}
\def\labelenumi{(\alph{enumi})}
\setcounter{enumi}{2}
\tightlist
\item
  考察练习 13 (b) 中的群 \(\mathbf{G}'\)，证明若将 \(\pi\) 换成 \(> \pi\) 的数，(b) 的结论不再成立。（在练习 13 的记号下，使用范数 \(ae_{1} + be_{2} + ce_{3} \mapsto (a^{2} + b^{2} + c^{2})^{\frac{1}{2}}\)，定义在 \(\mathrm{L}(\mathbf{G}')\) 上。）
\end{enumerate}

\begin{enumerate}
\def\labelenumi{\arabic{enumi}.}
\setcounter{enumi}{18}
\item
  设 G 是局部紧群，H 是 G 的闭子群。假设 H 是有限维单连通可解李群且 G/H 紧。证明 G 存在一个紧子群 L，使得 G 是 L 与 H 的半直积。（对 dim H 归纳。考虑 H 的最后一个非平凡导出群，并使用《积分理论》第 VII 章 § 3 命题 3。）
\item
  设 \(G\) 是有限维连通可解实李群。引入 §6 第 10 号命题 20 的证明中所用的记号 \(S, S', F, \sigma\)。
\end{enumerate}

\begin{enumerate}
\def\labelenumi{(\alph{enumi})}
\item
  利用命题 21，证明 \(\sigma\) 是从 S 到 S' 的底层实李群的一个李子群的同构。
\item
  推出 \(\tilde{G}\)（即 \(G\) 的万有复化）可识别为 \(S' / \sigma(F)\)，且从 \(G\) 到 \(\tilde{G}\) 的典范映射，是将 \(G\) 同构到 \(\tilde{G}\) 的底层实李群的一个李子群上的映射。
\end{enumerate}

¶ 21. (a) 设 G 是满足下列性质的单连通可解实李群：(α) L(G) 为 n 维，并具有一个 n - 1 维交换理想，该理想对应于 G 的子群 A；(β) G 的中心中存在一个不属于 A 的元素 σ。证明存在元素 \(x\)，属于 \(\mathrm{L}(\mathbf{G})\)，使 \(\exp x = \sigma\)。证明 \(\mathrm{L}(\mathbf{G})\) 是一个交换理想与一个具有如下基的理想的乘积：

\[
(x, a _ {1}, b _ {1}, \dots , a _ {k}, b _ {k})
\]

 其中 \(a_1, b_1, \ldots, a_k, b_k\) 属于 \(\mathrm{L}(\mathrm{A})\)，\([x, a_i] = 2\pi n_i b_i\)，\([x, b_i] = -2\pi n_i a_i\)（\(n_i \in \mathbf{Z} - \{0\}\)）对所有 \(i\) 成立。将练习 13 (d) 的结果推广到 G。

\begin{enumerate}
\def\labelenumi{(\alph{enumi})}
\setcounter{enumi}{1}
\tightlist
\item
  设 G 是有限维单连通可解实李群。设 D 是 G 的中心的离散子群。存在 L(G) 的一个基 \((x_{1}, x_{2}, \ldots, x_{n})\) 及一个整数 \(r \leqslant n\)，具有下列性质：(α) G 的每个元素都能唯一写成如下形式：
\end{enumerate}

\[
\left(\exp t _ {1} x _ {1}\right) \dots \left(\exp t _ {n} x _ {n}\right)
\]

  其中 \(t_1, \ldots, t_n\) 属于 \(\mathbf{R}\)；\((\beta)x_1, \ldots, x_r\) 两两可交换，且

\[
(\exp x _ {1}, \dots , \exp x _ {r})
\]

  是交换群 D 的一个基。（对 G 的维数归纳。设 a 是 L(G) 的极大交换理想，A 是相应的积分子群。于是 A 在 G 中闭，且 DA/A 是 G/A 的中心的离散子群，可以对其应用归纳假设，从而存在 L(G/A) 中的元素 \(x_1^*, \ldots, x_m^*\) 及一个整数 \(s\)。对于 \(1 \leqslant i \leqslant s\)，设 \(\sigma_i\) 是 D 中一个模 A 的类为 exp \(x_1^*\) 的元素。逐个使用 (a)，构造代表元 \(x_1, \ldots, x_m\)，使其成为 \(x_1^*, \ldots, x_m^*\) 的代表元，且 exp \(x_i = \sigma_i\)，以及 \([x_1, x_j] = 0\) 对 \(1 \leqslant i, j \leqslant s\) 成立。

\begin{enumerate}
\def\labelenumi{(\alph{enumi})}
\setcounter{enumi}{2}
\tightlist
\item
  由 (b) 推出每个有限维连通可解实李群都同态于空间 \(\mathbf{R}^n\times \mathbf{T}^m\)（\(m,n\) 为 \(\geqslant 0\) 的整数）。
\end{enumerate}

\begin{enumerate}
\def\labelenumi{\arabic{enumi}.}
\setcounter{enumi}{21}
\item
  设 \(G\) 是有限维连通实李群，且 \(L(G)\) 是约化的。则 \(G = (V \times S)/N\)，其中 \(V\) 是有限维实向量空间，其中 \(S\) 是单连通半单实李群，且 \(N\) 是 \(V \times S\) 的中心离散子群。于是 \(G/\overline{D^1}G\) 同构于 \(V/\mathrm{pr}_1\overline{N}\)。因此 \(G/\overline{D^1}N\) 紧，当且仅当 \(\mathrm{pr}_1\mathrm{N}\) 生成向量空间 \(V\)。
\item
  \begin{enumerate}
  \def\labelenumii{(\alph{enumii})}
  \tightlist
  \item
    设 G 是只有有限个连通分支的有限维交换实李群。G 紧的充要条件是：G 在复向量空间上的每个有限维解析线性表示都是半单的。（使用命题 32 的证明。）
  \end{enumerate}
\end{enumerate}

\begin{enumerate}
\def\labelenumi{(\alph{enumi})}
\setcounter{enumi}{1}
\tightlist
\item
  设 G 是只有有限个连通分支的有限维交换复李群。设 N 是 \(\exp_{G}\) 的核。
\end{enumerate}

N 在 C 上生成向量空间 \(L(G)\) 的充要条件是：G 的每个有限维解析线性表示都是半单的。（使用 (a)、引理 1 和命题 32。）

\begin{enumerate}
\def\labelenumi{\arabic{enumi}.}
\setcounter{enumi}{23}
\item
  李群 \(\mathbf{SL}(2,\mathbf{R})\) 是连通的且几乎单纯，但 \(\{I,-I\}\) 是 \(\mathbf{SL}(2,\mathbf{R})\) 的正规交换子群。
\item
  设 G 是几乎单纯的连通实或复李群，A 是 G 的正规子群。若 \(A \neq G\)，则 A 离散且为中心子群。（使用 §4 的练习 8。）因此，G 对其中心的商作为抽象群是单纯的。
\item
 设 G 是有限维连通实李群。假设 G 具有有限维单射连续线性表示 \(\rho\)。则 (G, G) 在 G 中闭。（设 R 是 G 的根基，S 是 G 的一个极大半单积分子群。利用练习 7 (b)，将问题化归为 G 是 S 与 R 的半直积的情形。根据第 I 章 § 6 命题 6，\(\rho\) 在 (G, R) 上是幺幂的。因此 \(\rho((G, R))\) 在线性群中闭，从而 (G, R) 在 G 中闭。）
\item
  设 G 是有限维可解单连通实李群，N 是 G 的最大连通幂零正规子群。则 G 具有一个在 N 上为幺幂的有限维单射连续线性表示。（对 G 的维数归纳。使用命题 20 以及第 I 章 § 7 定理 1。）†
\end{enumerate}

¶ 28. (a) 设 k 是特征为 0 的交换域，L 是 k 上的 n 维李代数，\(L_{1}\) 是一个 \((n-1)\) 维子代数，不含 L 的非零理想，且 \(a_{0}\) 是 L 中不属于 \(L_{1}\) 的元素。对于 \(i=2,3,\ldots\)，递归定义 \(L_{i}\) 为所有 \(x\in L_{i-1}\) 且满足 \([x,a_{0}]\in L_{i-1}\) 的元素组成的集合；记 \(L_{i}=L\) 对于 \(i\leqslant0\)。证明 \([L_{0},L_{i}]\subset L_{i+j}\)（对 \(i+j\) 归纳），然后证明对于 \(0\leqslant i\leqslant n\)，\(L_{i}\) 是 L 的余维为 i 的子代数（对 i 归纳）。

\begin{enumerate}
\def\labelenumi{(\alph{enumi})}
\setcounter{enumi}{1}
\tightlist
\item
  对于 0 \textless{} i \textless{} n，在 \(L_{i}\) 中选取元素 \(a_{i}\)，使其不属于 \(L_{i+1}\) 且满足 \([a_{0}, a_{i}] \equiv ia_{i-1} \pmod{L_{i}}\)。证明，对按字典序排列的有序对 \((i, j)\) 归纳，对于 \(0 \leqslant i < j < n\)、\(i + j - 1 < n\)，有
\end{enumerate}

\[
[ a _ {i}, a _ {j} ] \equiv (j - i) a _ {i + j - 1} (\mathrm{mod}. \mathrm{L} _ {i + j}).
\]

\begin{enumerate}
\def\labelenumi{(\alph{enumi})}
\setcounter{enumi}{2}
\item
  推出 L 或为 1 维，或为 2 维且非交换，或同构于 \(\mathfrak{sl}(2,k)\)。
\item
  设 A 是有限维实李代数。若 A 的子代数 B 存在满足 \(B_{1} \supset B\) 且 dim \(B_{1} = \dim B + 1\) 的子代数，则称 B 可扩张。记 \(R'(A)\) 为 A 的所有不可扩张子代数的交。记 \(R(A)\) 为 A 的最大理想，使其作为 A-模具有一个所有商模维数均为 1 的合成列。
\end{enumerate}

证明 \(\mathbf{R}'(\mathbf{A})\) 是 \(\mathbf{A}\) 的特征理想。（注意 \(\mathbf{R}'(\mathbf{A})\) 在 \(\operatorname{Aut}(\mathbf{A})\) 下稳定。）

\begin{enumerate}
\def\labelenumi{(\alph{enumi})}
\setcounter{enumi}{4}
\item
  证明 \(\mathbf{R}'(\mathbf{A})\supset \mathbf{R}(\mathbf{A})\)，且 \(\mathrm{R}'(\mathrm{A} / \mathrm{R}(\mathrm{A})) = \mathrm{R}'(\mathrm{A}) / \mathrm{R}(\mathrm{A})\)。
\item
  证明若 B 是 A 的子代数，则 R'(B) ⊃ R'(A) ∩ B。
\item
  若 A 可解，则 \(\mathrm{R}^{\prime}(\mathrm{A}) = \mathrm{R}(\mathrm{A})\)。（根据 (e)，可假设 \(\mathrm{R}(\mathrm{A}) = \{0\}\)。设 \(\mathrm{R}^{\prime}(\mathrm{A}) \neq \{0\}\)。由 (d)，存在包含于 \(\mathrm{R}^{\prime}(\mathrm{A})\) 中的 A 的极小非零理想 I。利用 (f) 及第 I 章 §5 定理 1 的推论 1，证明 dim I = 1，从而得到矛盾。）
\item
  设 P 是 \(\mathrm{R}^{\prime}(\mathrm{A})\) 的根基，B 是 A 的半单子代数，且 \(C = P + B\) 根据 (d) 是 A 的子代数。利用 (f)，证明相对于 P 的适当基，\(ad_{P}B\) 可表示为三角形。推出 \([P, B] = \{0\}\)。
\end{enumerate}

\((k)\) R(A) 是 \(\mathrm{R}'(\mathrm{A})\) 的根基。（使用 \((e)\)、\((f)\)、\((g)\)、\((h)\) 以及 A 的 Levi 分解。）设 \(\mathrm{R}'(\mathrm{A}) = \mathrm{R}(\mathrm{A}) \oplus \mathrm{T}\) 是 \(\mathrm{R}'(\mathrm{A})\) 的 Levi 分解。根据 \((h)\)，\(\mathrm{R}'(\mathrm{A}) = \mathrm{R}(\mathrm{A}) \times \mathrm{T}\)。于是 \(\mathrm{R}(\mathrm{T}) = \{0\}\)。将 \((c)\) 应用于 T 的单因子，推出 \(\mathrm{T} = \{0\}\)。因此已经证明 \(\mathrm{R}(\mathrm{A}) = \mathrm{R}'(\mathrm{A})\)。

\begin{enumerate}
\def\labelenumi{(\alph{enumi})}
\setcounter{enumi}{11}
\item
  设 G 是李代数为 A 的连通实李群。令 \(\mathcal{S}(G)\) 为 G 的子群 N 的集合，其中存在具有下列性质的递减子群序列 \((\mathrm{N}_m, \mathrm{N}_{m-1}, \ldots, \mathrm{N}_0)\)：\(\mathrm{N}_m = \mathrm{N}\)，\(\mathrm{N}_0 = \{\varepsilon\}\)，每个 \(\mathrm{N}_i\) 都是 G 的正规连通李子群，且 \(\dim \mathrm{N}_i / \mathrm{N}_{i-1} = 1\) 对所有 \(i > 0\) 成立。证明，李代数为 R(A) 的 G 的积分子群 R(G) 是 \(\mathcal{S}(G)\) 的最大元。（对 \(\dim \mathrm{R(A)}\) 归纳。设 I 是 A 的 1 维理想，N 是 G 中相应的积分子群。对 N 取商。若 N 不闭，使用 § 6 的练习 21 (a)。）
\item
  若 G 的李子群 H 存在满足 \(H_{1} \supset H\) 且 \(\dim H_{1} = \dim H + 1\) 的李子群，则称 H 可扩张。记 \(R'(G)\) 为 G 的所有不可扩张连通李子群的交。
\end{enumerate}

设 B 是 A 的不可扩张子代数。证明 G 中相应的积分子群是闭的。（使用命题 5。）推出

\[
\mathrm{L} (\mathrm{R} ^ {\prime} (\mathrm{G})) \subset \mathrm{R} (\mathrm{A}).
\]

从而 \(\mathrm{R}^{\prime}(\mathrm{G})\subset\mathrm{R}(\mathrm{G})\)。

\begin{enumerate}
\def\labelenumi{(\alph{enumi})}
\setcounter{enumi}{13}
\tightlist
\item
  证明 \(\mathrm{R}(\mathrm{G}) = \mathrm{R}'(\mathrm{G})\)。（将其化归为 \(\mathrm{R}'(\mathrm{G}) = \{e\}\) 的情形。然后使用 §6 的练习 22。）†
\end{enumerate}

\begin{enumerate}
\def\labelenumi{\arabic{enumi}.}
\setcounter{enumi}{28}
\tightlist
\item
  实李群 G 称为 (N) 型，如果它是有限维的，
\end{enumerate}

幂零、连通且单连通。若 G 是这样的群且 g 是其李代数，则当赋予 g 以由豪斯多夫法则定义的群结构时，映射 exp: g → G 是同构（参见第 II 章 §6 第 5 号注记 3）。记 log: G → g 为逆同构。

\begin{enumerate}
\def\labelenumi{(\alph{enumi})}
\item
  设 V 是 g 的一个向量 Q-子空间。证明下列条件等价：
\item
  V 是 \(\mathbf{Q}\) 上的 \(\mathfrak{g}\) 的李子代数；
\end{enumerate}

\begin{enumerate}
\def\labelenumi{(\roman{enumi})}
\setcounter{enumi}{1}
\tightlist
\item
  \(\exp(V)\) 是 G 的子群。
\end{enumerate}

（使用第 II 章 § 6 的练习 5。）

G 的子群 H 要具有 \(\exp(V)\) 的形式（其中 V 是 g 的向量 Q-子空间），充要条件是 H 在 G 中饱和（第 II 章 §4，练习 14），即关系 \(x \in G\)、\(x^{n} \in H\)、\(n \neq 0\) 蕴含 \(x \in H\)（同上）。若此条件满足，证明 H 是 G 的积分子群，当且仅当 \(\log(H)\) 是 g 的李 R-子代数。

\begin{enumerate}
\def\labelenumi{(\alph{enumi})}
\setcounter{enumi}{1}
\item
  设 V 是 g 的一个有限维李 Q-子代数，\((e_{1},\ldots,e_{m})\) 是 V 在 Q 上的一个基，\(\Lambda\) 是由 \(e_{1},\ldots,e_{m}\) 生成的 V 的子群。利用豪斯多夫法则是多项式这一事实，证明存在一个整数 \(d\geqslant1\)，使得对于 d 的每个非零倍数整数 r，\(\exp(r\Lambda)\) 都是 G 的子群。设 d 为这样的整数，r 为 d 的倍数。证明 \(\exp(r\Lambda)\) 离散，当且仅当 \((e_{1},\ldots,e_{m})\) 是 R 上的自由族，即 V \(\otimes_{q}R\) 到 g 的典范映射是单射。假设此条件成立；证明 G/exp(r \(\Lambda\) ) 紧，当且仅当 V \(\otimes_{q}R\to g\) 是双射，即 V 是 g 的一个 Q-型。（为证明充分性，对 g 的幂零类归纳。）
\item
  反之，设 \(\Gamma\) 是 G 的离散子群，令 \(\Gamma\) 为其在 G 中的饱和（第 II 章，同上），并令 \(\mathfrak{g}_{\Gamma} = \log (\Gamma) = \mathbf{Q}. \log (\Gamma)\) 为相应的李 \(\mathbf{Q}\)-子代数。假设 \(\mathbf{R}, \mathfrak{g} = \mathfrak{g}\)，即 \(\Gamma\) 不包含于 G 的任何不同的积分子群中。设 \(z\) 是满足 \(\neq 1\) 的 \(\Gamma\) 的中心元素，令 \(x = \log (z)\)。证明 \(x\) 属于 \(\mathfrak{g}\) 的中心。若 \(\mathrm{X} = \exp (\mathbf{R}x)\)，证明 \(\mathrm{X} / (\Gamma \cap \mathrm{X})\) 是紧的，且 \(\Gamma\) 在 \(\mathrm{G} / \mathrm{X}\) 中的像是离散的。通过对 dim G 归纳，推出 \(\mathrm{G} / \Gamma\) 紧，且 \(\mathfrak{g}_{\Gamma}\) 是 \(\mathfrak{g}\) 的 G-型。若 \(\Lambda\) 是 \(\mathfrak{g}_{\Gamma}\) 的格，证明存在整数 \(d \neq 0\)，使得 \(\Gamma\) 包含 \(\exp (d\Lambda)\)，并且 \(\exp (d\Lambda)\) 在 \(\Gamma\) 中的指数有限。
\end{enumerate}

设 H 是李代数为 \(\mathfrak{h}\) 的 G 的积分子群。证明 \(\mathrm{H} / (\mathrm{H} \cap \Gamma)\) 紧，当且仅当 \(\mathfrak{h}\) 关于 \(\mathbf{Q}\)-结构 \(g_{\Gamma}\) 是有理的（特别地，当 \(\mathfrak{h}\) 是 \(g\) 的下中心列或上中心列中的一项时成立）。

\begin{enumerate}
\def\labelenumi{(\alph{enumi})}
\setcounter{enumi}{3}
\item
  设 \(\Gamma\) 是 \(G\) 的离散子群，令 \(H\) 是 \(G\) 的最小积分子群，包含 \(\Gamma\)，且 \(\mathfrak{h}\) 是其李代数。证明 \(H / \Gamma\) 紧，且 \(\mathfrak{g} \otimes_{\mathbf{q}} \mathbf{R} \to \mathfrak{g}\) 是单射，其像为 \(\mathfrak{h}\)。（将 (c) 应用于幂零群 \(H\)。）
\item
  设 \(\Gamma\) 是 \(G\) 的离散子群。证明存在基 \((x_{1},\ldots ,x_{q})\) 作为 \(L(G)\) 的基，具有下列性质：
\item
  对所有 \(i \in [1, q]\)，\(\mathbf{R}x_i + \cdots + \mathbf{R}x_q\) 是记为 \(\mathfrak{n}_i\) 的 \(\mathbf{R}x_{i-1} + \cdots + \mathbf{R}x_q\) 的理想；
\end{enumerate}

\begin{enumerate}
\def\labelenumi{(\roman{enumi})}
\setcounter{enumi}{1}
\tightlist
\item
  存在 \(p \in (1, q]\)，使得 \(\Gamma\) 恰为如下乘积的集合
\end{enumerate}

\[
\exp (m _ {p} x _ {p}) \exp (m _ {p + 1} x _ {p + 1}) \dots \exp (m _ {q} x _ {q})
\]

其中 \(m_{p}, \ldots, m_{q}\) 属于 Z；

\begin{enumerate}
\def\labelenumi{(\roman{enumi})}
\setcounter{enumi}{2}
\tightlist
\item
  若 \(G / \Gamma\) 紧，则 \(\{n_1, n_2, \ldots, n_q\}\) 包含 \(g\) 的下中心列。（利用 (d) 和命题 16，将其化归为 \(G / \Gamma\) 紧的情形。然后利用 (c) 对 \(\dim G\) 归纳。）
\end{enumerate}

\begin{enumerate}
\def\labelenumi{(\arabic{enumi})}
\setcounter{enumi}{29}
\tightlist
\item
  举出一个 7 维实幂零李代数的例子，使其不存在一个基，关于该基的结构常数为有理数（参见第 I 章 § 4 的练习 18）。推出相应的 (N) 型群（练习 29）没有带紧子群的离散子群。（使用练习 29。）
\end{enumerate}

\begin{enumerate}
\def\labelenumi{\arabic{enumi}.}
\setcounter{enumi}{30}
\item
  设 G 和 \(G'\) 是两个 (N) 型李群（练习 29），\(\Gamma\) 是 G 的离散子群且 \(G/\Gamma\) 紧。证明每个同态 \(f: \Gamma \to G'\) 都能唯一延拓为从 G 到 \(G'\) 的李群态射（首先将 f 延拓到 \(\tilde{\Gamma}\) 的饱和 \(\Gamma\)，并导出从李 Q-代数 \(\log(\tilde{\Gamma})\) 到 \(G'\) 的李代数的一个同态；参见练习 29）。
\item
  设 G 是 (N) 型李群（练习 29），\(\Gamma\) 是 G 的离散子群。证明下列条件等价：
\end{enumerate}

\begin{enumerate}
\def\labelenumi{(\alph{enumi})}
\item
  G/Γ 是紧的；
\item
  \(\mathrm{G} / \Gamma\) 的测度有限（相对于一个非零的 G-不变正测度）；
\item
  每个包含 \(\Gamma\) 的 G 的积分子群都等于 G。（使用练习 29。）
\end{enumerate}

¶ 33. 设 \(\Gamma\) 是一个群。证明下列条件等价：(a) \(\Gamma\) 是幂零、无挠且有限生成的；

\begin{enumerate}
\def\labelenumi{(\alph{enumi})}
\setcounter{enumi}{1}
\item
  存在一个包含 \(\Gamma\) 作为离散子群的 (N) 型李群（练习 29）；
\item
  存在一个 (N) 型李群 G（练习 29），包含 \(\Gamma\) 作为离散子群，且 \(G / \Gamma\) 紧。
\end{enumerate}

（\(\langle b\rangle\) 与 \(\langle c\rangle\) 的等价性由练习 29 得出。蕴含 \((\varepsilon)\Rightarrow \langle a\rangle\) 通过对 \(\dim (\mathrm{G})\) 归纳并使用练习 29 \((c)\) 的方法得到证明。为证明 \((a)\Rightarrow (c)\)，先证明该李 \(\mathbf{Q}\)-代数（附属于 \(\Gamma\) 的饱和）是有限维的（参见第 II 章 § 6 的练习 4）；然后取该代数与 \(\mathbf{R}\) 的张量积以及相应的 (N) 型群。）

\begin{enumerate}
\def\labelenumi{\arabic{enumi}.}
\setcounter{enumi}{33}
\tightlist
\item
  设 \(G\) 是有限维连通实或复李群，\(\rho\) 是 \(G\) 在有限维复向量空间 \(V\) 上的解析线性表示。假设 \(\rho\) 是半单的。群 \(G\) 通过自同构作用于代数 \(S(V^*)\)，即 \(V\) 上的多项式函数代数。
\end{enumerate}

\begin{enumerate}
\def\labelenumi{(\alph{enumi})}
\item
  设 \(\mathsf{S}(\mathbf{V}^{*})^{\mathbb{G}}\) 是 \(\mathsf{S}(\mathbf{V}^{*})\) 的 G-不变元素的集合。存在投影算子 \(p\)，从 \(S(V^{*})\) 到 \(S(V^{*})^{\mathbb{G}}\)，它与 \(G\) 的运算交换，并使 \(G\)-稳定的 \(S(V^{*})\) 的每个向量子空间保持稳定。
\item
  设 I、J 是 \(S(V^*)\) 的在 G 下稳定的理想。令 A（分别为 B）是 I（分别为 J）在 V 中的零点集。假设 \(A \cap B = \varnothing\)。于是存在 \(u \in I\)，使 \(u\) 是 G-不变的，且在 B 上 \(u = 1\)。（根据《交换代数》第 7 章 §3 第 3 号命题 2，存在 \(v \in I\)、\(w \in J\)，使 \(v + w = 1\)。令 \(u = pv\)。证明 \(u\) 具有所需性质。）
\end{enumerate}

\begin{enumerate}
\def\labelenumi{\arabic{enumi}.}
\setcounter{enumi}{34}
\tightlist
\item
  设 \(\mathbf{G}\) 是 \(\mathbf{GL}(n,\mathbf{R})\) 的紧子群。
\end{enumerate}

\begin{enumerate}
\def\labelenumi{(\alph{enumi})}
\item
  设 A、B 是 \(\mathbf{R}^{\mathbf{n}}\) 中不交的紧 G-不变子集。存在 G-不变的多项式函数 \(u\)，定义在 \(\mathbf{R}^{\mathbf{n}}\) 上，使得 \(u = 1\) 在 \(\mathrm{A}\) 上成立且 \(u = 0\) 在 B 上成立。（存在连续实值函数 \(v\)，定义在 \(\mathbf{R}^{\mathbf{n}}\) 上，使得 \(v = 1\) 在 A 上成立且 \(v = 0\) 在 B 上成立。由 Stone-Weierstrass 定理，存在多项式函数 \(w\)，定义在 \(\mathbf{R}^{\mathbf{n}}\) 上，使得 \(|v - w| \leqslant \frac{1}{3}\) 在 \(\mathrm{A} \cup \mathrm{B}\) 上成立。通过相对于 G 的归一化 Haar 测度的积分过程得到 \(u\)。）
\item
  设 \(x \in \mathbf{R}^n\)。则 \(Gx\) 是 \(\mathbf{R}^n\) 中有限个 G-不变多项式的零点集。（使用 (a) 以及第 8 号定理 1 的推论 2。）
\end{enumerate}

\begin{enumerate}
\def\labelenumi{\arabic{enumi}.}
\setcounter{enumi}{35}
\tightlist
\item
  在 \(\mathbf{SL}(2, \mathbf{R})\) 中，能写成 \((a, b)\) 形式的元素（其中 \(a, b\) 属于 \(\mathbf{SL}(2, \mathbf{R})\)）正是满足 \(\neq -1\) 的元素。
\end{enumerate}

¶ 37. 设 G 是紧实李群，G' 是有限维连通实李群。假设 L(G) 是单纯的。设 \(\rho: G \to G'\) 是抽象群同态。假设 G 中存在 \(e_{\alpha}\) 的邻域 V，使得 \(\rho(V)\) 相对紧。则 \(\rho\) 连续。（设 V' 是 G' 中 \(e_{\alpha'}\) 的邻域。设 V'\,' \(\subset\) V' 是 \(e_{\alpha'}\) 的邻域，使得

\[
(x \in \mathrm{V} ^ {\prime \prime}, x _ {j} \in \rho (\mathrm{V}), y _ {j} \in \rho (\mathrm{V}) \text {   for   } 1 \leqslant j \leqslant n) \Rightarrow \left(\prod_ {i = 1} ^ {n} x _ {i} (y _ {i}, x) x _ {i} ^ {- 1} \in \mathrm{V} ^ {\prime}\right)
\]

其中 \(n = \dim G\)。可假设 \(\rho\) 非平凡。由练习 25，\(\operatorname{Ker} \rho\) 是有限的。因此 \(\rho(G)\) 不可数，因而不是离散的。于是存在 \(g \in G\)，其中心化子不开放，且 \(\rho(g) \in V''\)。使用 §4 练习 9 及其记号，有 \(\rho(M(g, n, V)) \subset V'\)，且 \(M(g, n, V)\) 是 G 中 \(e_G\) 的邻域。）

\begin{enumerate}
\def\labelenumi{\arabic{enumi}.}
\setcounter{enumi}{37}
\tightlist
\item
  设 \(G\) 是有限维实李群，\(G_0\) 是其单位元连通分支。考虑以下性质：
\end{enumerate}

\begin{enumerate}
\def\labelenumi{(\Alph{enumi})}
\setcounter{enumi}{5}
\tightlist
\item
  Ad(G) 在 Aut L(G) 中闭。
\end{enumerate}

证明在下列每种情形中 G 都具有性质 (F)：

\begin{enumerate}
\def\labelenumi{(\roman{enumi})}
\item
  G 连通且幂零；
\item
  L(G) 的每个导子都是内导子；
\item
  \(\mathrm{G}_0\) 具有性质 (F)，且 \(\mathrm{G} / \mathrm{G}_0\) 是有限的；
\item
  G 是上三角群；
\item
  \(G_{0}\) 是半单的。
\end{enumerate}

\begin{enumerate}
\def\labelenumi{\arabic{enumi}.}
\setcounter{enumi}{38}
\tightlist
\item
  设 G 是有限维连通实李群，H 是 G 的积分子群，\(\bar{H}\) 是 H 在 G 中的闭包。假设 H 具有性质 (F)（练习 38）。
\end{enumerate}

\begin{enumerate}
\def\labelenumi{(\alph{enumi})}
\item
  设 \(x \in \overline{\mathbf{H}}\)。则 \(\mathrm{Ad}_{\mathrm{L(G)}}x\) 保持 \(\mathrm{L(H)}\) 稳定（第 2 号，命题 5）；令 \(u(x)\) 为其在 \(\mathrm{L(H)}\) 上的限制。于是 \(u(\overline{\mathrm{H}}) = \mathrm{Ad}_{\mathrm{L(H)}}(\mathrm{H})\)。（注意 \(\mathrm{Ad}_{\mathrm{L(H)}}(\mathrm{H})\) 在 \(u(\overline{\mathrm{H}})\) 中稠密，并应用性质 (F)。）
\item
  设 C 是 \(\overline{\mathrm{H}}\) 的中心。则 \(\overline{\mathrm{H}} = \mathrm{C.H}\)，且 C 是 H 的中心的闭包。（根据 (a)，若 \(x \in \overline{\mathrm{H}}\)，则存在 \(y \in \mathrm{H}\)，使 \(\mathrm{Ad}_{\mathrm{L(H)}} x = \mathrm{Ad}_{\mathrm{L(H)}} y\)。于是 \(\mathrm{Ad}(x^{-1}y)\) 在 L(H) 上是恒等的，从而 \(x^{-1}y \in Z_{\overline{\mathrm{H}}}(\mathrm{H})\)，且 \(x \in Z_{\overline{\mathrm{H}}}(\mathrm{H})\)。H。由连续性，\(Z_{\overline{\mathrm{H}}}(\mathrm{H})\) 等于 C。群 C 在 \(\overline{\mathrm{H}}\) 中闭，因而在 G 中也闭；所以 \(\overline{\mathrm{C}} \cap \overline{\mathrm{H}} \subset \mathrm{C}\)。设 \(x \in \mathrm{C}\)。存在 H 中的序列 \((x_n)\)，使 \(x_n\) 趋于 \(x\)。于是 \(\mathrm{Ad}_{\mathrm{L(H)}} x_n\) 趋于 1，从而存在 H 中的序列 \((y_n)\)，使 \(y_n\) 趋于 \(e\)，且 \(\mathrm{Ad}_{\mathrm{L(H)}} y_n = \mathrm{Ad}_{\mathrm{L(H)}} x_n\)。于是 \(y_n^{-1}x_n \in \mathrm{C} \cap \mathrm{H}\)，且 \(y_n^{-1}x_n\) 趋于 \(x\)。）
\item
  \(\mathbf{H} = \overline{\mathbf{H}}\) 当且仅当 \(\mathbf{H}\) 的中心在 \(\mathbf{G}\) 中闭。（使用 (b)。）
\end{enumerate}

\begin{enumerate}
\def\labelenumi{\arabic{enumi}.}
\setcounter{enumi}{39}
\tightlist
\item
  设 \(\mathbf{H}\) 是有限维实李群，\(\mathbf{H}_0\) 是其单位元连通分支。
\end{enumerate}

\begin{enumerate}
\def\labelenumi{(\alph{enumi})}
\item
  假设 \(H_{0}\) 满足练习 38 的条件 (F)，\(H/H_{0}\) 有限，且 \(H_{0}\) 的中心紧。设 G 是有限维实李群，\(f:H\to G\) 是具有离散核的连续同态。则 \(f(\mathrm{H})\) 在 G 中闭。（将其化归为 H 连通的情形。由于 f 的核是离散的，\(f(\mathrm{H})\) 的中心是 H 的中心在 f 下的像（引理 1），因而在 G 中闭。应用练习 39 (c)。）
\item
  假设 \(H_{0}\) 是半单的且 \(H/H_{0}\) 有限。则 H 在有限维连续线性表示下的像在线性群中闭。（应用 (a) 和练习 7 (b)。）
\end{enumerate}

\begin{enumerate}
\def\labelenumi{\arabic{enumi}.}
\setcounter{enumi}{40}
\tightlist
\item
  设 \(G\) 是有限维连通实李群，\(\rho: G \to \mathbf{GL}(n, \mathbf{C})\) 是具有有限核的连续同态。则 \(\rho((G, G))\) 在 \(\mathbf{GL}(n, \mathbf{C})\) 中闭。（利用练习 7 (b) 和 40 (a)，将其化归为 \(G\) 是半单群 \(S\) 与其根基 \(R\) 的半直积的情形。\(\rho((G, R))\) 的每个元素都是幺幂的，故 \(\rho((G, R))\) 闭。向量子空间 \(V_1\) 由 \(\rho((G, R))\) 的不动点组成，满足 \(\neq \{0\}\)。它在 \(\rho(G)\) 下稳定。考虑 \(\mathbf{C}^n / V_1\)，对 \(n\) 归纳并使用 \(\rho(S)\) 的完全可约性，可以看出能够写成
\end{enumerate}

\[
\mathbf {C} ^ {n} = \mathrm{V} _ {1} \oplus \dots \oplus \mathrm{V} _ {q}
\]

其中每个 \(V_{i}\) 在 \(\rho(S)\) 下稳定，且对所有 i 有 \(\rho((G,R))(V_{i})\subset V_{1}\oplus\cdots\oplus V_{i-1}\)。另一方面，\(\rho(S)\) 是闭的（练习 40 (b)）。最后，\(\rho((G,G))=\rho(S)\cdot\rho((G,R)).\)

\begin{enumerate}
\def\labelenumi{\arabic{enumi}.}
\setcounter{enumi}{41}
\tightlist
\item
  设 \(G\) 是有限维连通实李群。假设
\end{enumerate}

G 具有一个单射连续线性表示 \(\rho: \mathrm{G} \to \mathbf{GL}(n, \mathbf{C})\)。则 G 具有一个线性表示 \(\sigma: \mathrm{G} \to \mathbf{GL}(n', \mathbf{C})\)，它把 G 同态到 \(\mathbf{GL}(n', \mathbf{C})\) 的一个闭子群上。（根据练习 41，\(\rho((\mathrm{G}, \mathrm{G}))\) 在 \(\mathbf{GL}(n, \mathbf{C})\) 中闭，因而 (G, G) 在 G 中闭。令 \(\rho: \mathrm{G} \to (\mathrm{G}, \mathrm{G})\) 为典范态射，\(\tau\) 为一个像闭的单射线性表示，其像 \(\mathrm{G}/(\mathrm{G}, \mathrm{G})\) 连通且交换。取 \(\sigma = \rho \oplus (\tau \circ p)\)。）

